\documentclass[a4paper]{article}
\usepackage{amsfonts}
\usepackage{amsmath}
\usepackage{amssymb}
\usepackage{graphicx}  
\usepackage{cancel}

\newcommand{\degree}{^{\circ}}
\newcommand{\cis}{\operatorname{cis}}
\newcommand{\Bgsin}{\operatorname{Bgsin}}
\newcommand{\Bgcos}{\operatorname{Bgcos}}
\newcommand{\Bgtan}{\operatorname{Bgtan}}
\newcommand{\limit}{\lim\limits}

\begin{document}
  
\noindent \large Auteur: Rune Suy \\
\noindent \large Studierichting: Informatica\\
\noindent \large Oefeningenreeks 2E nr. 26.23\\

\medskip

\normalsize

Los de volgende logaritmische vergelijking op: \\

$\log_{x+3}x \cdot \log_{10}x + \log_{10}x - \log_{x+3}x = 0$\\

$\Leftrightarrow \log_{x+3} x \cdot \log_{10}x = \log_{x+3}x - \log_{10}x$\\

$\Leftrightarrow \log_{10}x = \dfrac{\log_{x+3}x - \log_{10}x}{\log_{x+3}} = 1 - \dfrac{\log_{10}x}{\log_{x+3}x} = 1 - \dfrac{\log_{10}x}{\dfrac{\log_{10}x}{\log_{10}(x+3)}}$\\

Geval 1:  $x \neq 1$\\

$log_{10}x = 1 - \dfrac{\cancel{\log_{10}x}}{\dfrac{\cancel{\log_{10}x}}{\log_{10}(x+3)}} = 1 - \log_{10}(x+3)$\\

$\Leftrightarrow \log_{10}(x \cdot (x+3)) = 1$\\

$\Leftrightarrow 10 = x^2 + 3x$\\

$\Leftrightarrow x^2 + x - 10 = 0$\\

$D = 49 = 7^2$\\

$x_1 = \dfrac{-3-7}{2} = \cancel{-5}$\\

$x_2 = \dfrac{-3+7}{2} = 2$\\

Geval 2: $x=1$\\

$\log_{x+3}1 \cdot \log_{10}1 + \log_{10}1 - \log_{x+3}1 = 0$\\
$$ V = \{2, 1\}$$

\begin{thebibliography}{999}
%\bibitem{a} Referentie
\end{thebibliography}
\end{document} 
